\section{Vision: A Road-Network Congestion Signal}
\label{sec:vision}

The central claim is not that vehicles should form a fully cooperative swarm. It is that they should expose just enough local traffic state for nearby vehicles to make better independent decisions. This distinction matters for deployment. Command-style V2V coordination creates hard questions about authority, liability, and failure: who tells whom to slow down, which vehicle must occupy a lane, and what happens when messages are delayed? Self-regulating cars should avoid that contract. Each vehicle should remain the sole decision maker for its own motion, and network-level behavior should emerge from repeated local damping rather than from an agreed-upon plan.

\noindent\textbf{Semantic, aggregate exchange.}
We therefore define V2V communication as a bounded semantic API. A sender contributes compact traffic-state summaries, and a receiver aggregates whatever it heard into its own local observation. The API is intentionally lower level than a plan and higher level than raw sensing. It gives a vehicle observability beyond its camera horizon without turning the network into a centralized planner, and without publishing anything that identifies or tracks an individual vehicle. Concretely, the contract admits and forbids two disjoint classes of state:
\begin{itemize}[nosep,leftmargin=1.2em]
\item \textbf{Exchanged:} nearby vehicle counts and density for both AVs and non-AVs, mean speed, lane distribution, local queue estimate, downstream congestion estimate, segment target speed, and merge pressure.
\item \textbf{Never exchanged:} raw video, stable identities, direct acceleration or braking commands, joint lane plans, assignments of AVs to lanes, and coordinated roadblock behavior.
\end{itemize}

\noindent\textbf{The contract.}
More formally, vehicle \(i\) maintains a local peer set
\[
N_i(t)=\{j: \mathrm{dist}(i,j) \le R,\; t-\tau_j \le T\},
\]
where \(R\) is the aggregation radius, \(T\) is the freshness horizon, and \(\tau_j\) is the receive time of the latest summary from peer \(j\). The receiver reduces the summaries from \(N_i(t)\) to a single aggregate observation \(a_i(t)=A(\{b_j:j\in N_i(t)\})\), using count sums and confidence-weighted means over fields such as local count, mean speed, density, lane distribution, queue estimate, and merge pressure. The exact radio and aggregation operator can vary, but three semantics should not: expired peers are removed rather than extrapolated, unheard peers map to explicit neutral fields rather than zeros with hidden meaning, and \(a_i(t)\) is advisory evidence rather than a command. In our prototype interface, this substrate is realized as low-rate broadcasts at 2\,Hz, with a 2\,s peer time-to-live (TTL) and a 150\,m radius. This is a stand-in for a DSRC or C-V2X radio, not a claim about a specific one. If no peer is heard, the receiver falls back to neutral values and behaves as a local-only controller; if conflicting summaries arrive, it downweights low-confidence state rather than acting on any single message; and if bandwidth is scarce, it can trade rate or precision while still computing and filtering its own advisory.

\noindent\textbf{Receiver-side semantics and privacy.}
The receiver-side rules are as important as the message format. A self-regulating car treats V2V state as an observation hint, not as authority. The local digital twin records when an aggregate was last refreshed, whether a field came from onboard sensing or V2V, and whether it was neutralized because no peers were heard. This lets the policy distinguish "I see no cooperating vehicles" from "the radio is stale", and lets the safety layer reject advisories that depend on low-confidence state. It also keeps the privacy surface small \emph{by construction}: because the contract's unit is a neighborhood aggregate, a participating vehicle need not publish camera frames, a persistent cross-session identifier, or its own trajectory for its neighbors to benefit from the density and speed information they collectively induce. Realizing that guarantee end to end, so that no raw per-vehicle position ever leaves the sender, is an open problem we return to in Section~\ref{sec:discussion}.

\noindent\textbf{Why aggregate state is enough.}
Highway regulation does not require every vehicle to know the full network state, only early evidence that local flow is approaching a bad operating region: density rising near a merge, queues forming downstream, or speed variance increasing before a shockwave. This suffices because the control objective is physical and local to evaluate. The fundamental diagram relates flow to density, rising to a critical operating point and then collapsing as interactions dominate~\cite{greenshields1935study,lighthill1955kinematic,richards1956shock}; a regulator need only keep density below that threshold and then hold speed high and variance low, and both quantities are precisely what the aggregate signal conveys. A vehicle that damps its own speed variance when it sees rising local pressure removes energy from an incipient wave before it steepens, and when enough vehicles do this independently the stream is held near critical density with no shared plan. Onboard perception supplies immediate headway and leader and follower state, while aggregate V2V fills the blind spots beyond occlusions and around merges, so the observation is local but no longer myopic. The target is thus not consensus but local control under partial information, which is why the communication abstraction, rather than the optimizer, is the central design question.

\noindent\textbf{When is the signal active? Penetration and density.}
An important concern is that at low adoption a vehicle may hear no cooperating peers at all, leaving it a local-only controller. The saving structure is that the number of peers scales with the very density the signal exists to watch. Consider a receiver with radius \(R=150\)m on a three-lane highway, roughly a \(0.3\)km\,\(\times\,3\)-lane window. Near the critical density of about \(25\) veh/km/lane the window holds \(\sim\!22\) vehicles, of which a fraction \(p\) cooperate: about \(0.2\) peers at \(p=1\%\), \(1.1\) at \(5\%\), and \(2.3\) at \(10\%\). As flow degrades toward jam density (\(\sim\!100\) veh/km/lane) the same window holds \(\sim\!90\) vehicles, giving \(\sim\!4.5\) peers even at \(p=5\%\). The aggregate substrate therefore activates monotonically as local density rises toward the bad operating region it is meant to detect, and is quietest exactly when flow is healthy and sparse. We note that this implies a floor: below a few percent penetration near free flow a vehicle often hears nothing and must fall back to local control, and the fields that depend on beyond-horizon evidence, downstream congestion, queue estimate, and merge pressure, are the first to revert to neutral, while onboard headway and leader and follower state remain. Establishing where on this penetration--density curve the aggregate begins to change the network outcome is a central empirical question for the evaluation.

\noindent\textbf{Does decentralization give up too much?}
Independent local controllers, with no shared plan, might be expected to oscillate or to fall short of a coordinated optimum. Two things make us optimistic. First, the sparse-actuator result is that even a \emph{single} uncoordinated vehicle, with no V2V at all, already damps stop-and-go waves on both ring tracks and open roads~\cite{stern2018dissipation,lee2024traffic}; decentralized local action is thus not the speculative part, and the aggregate signal only lets each already-effective actuator act less myopically. Second, the property that must be protected as actuators are added is string stability, the requirement that disturbances not amplify as they propagate upstream through a chain of followers~\cite{swaroop1996string}; this gives a concrete, testable target for the safety layer and the evaluation. What the aggregate changes is the information each controller acts on, not the number of controllers, so the decentralization question reduces to whether better local observability improves a control law that is already known to help in the fully myopic case.

\noindent\textbf{Relationship to network control.}
The analogy to communication networks is direct: backpressure scheduling stabilizes queues by routing according to local pressure differences~\cite{tassiulas1990stability}, and self-regulating cars seek an analogous pressure signal for physical traffic. The difference is that packets do not brake hard when startled and humans do not obey schedulers, which is why the communication contract must be paired with an ego-centric model of the scene and a safety layer that can reject any unsafe or obstructive advisory.
